\documentclass[12pt]{article}

\usepackage[utf8]{inputenc}
\usepackage[T1]{fontenc}
\usepackage{chemfig}
\usepackage{geometry}

\geometry{margin=2.5cm}

\title{Ejemplos de Fórmulas Químicas con \texttt{chemfig}}
\author{}
\date{}

\begin{document}

\maketitle

\section{Introducción}

La librería \texttt{chemfig} permite representar moléculas y estructuras químicas de forma sencilla.

\subsection{Enlace simple}

\[
\chemfig{A-B}
\]

\subsection{Enlace doble}

\[
\chemfig{A=B}
\]

\section{Anillos}

\subsection{Anillo de cinco miembros}

\[
\chemfig{A*5(-B=C-D-E=)}
\]

\subsection{Anillo incompleto}

\[
\chemfig{A*5(-B=C-D)}
\]

\section{Molécula de Cafeína}

La siguiente estructura corresponde a una representación simplificada de la cafeína:

\[
\chemfig{*6((=O)-N(-)-(*5(-N=-N(-)-))=-(=O)-N(-)-)}
\]

\section{Acetaldehído}

También es posible personalizar el estilo de los enlaces.

\bigskip

{\Large
\setchemfig{
    atom sep=2em,
    bond style={
        line width=1pt,
        red,
        dash pattern=on 2pt off 2pt
    }
}

\chemname
    {\chemfig{H-C(-[2]H)(-[6]H)-C(=[1]O)-[7]H}}
    {Acetaldehído}
}

\section{Conclusión}

La librería \texttt{chemfig} facilita la creación de fórmulas químicas, estructuras cíclicas y moléculas complejas directamente en \LaTeX, permitiendo además personalizar la apariencia de los enlaces y etiquetas.

\end{document}
